\documentclass[10pt,a4paper]{amsart}
\usepackage[margin=19mm]{geometry}
\usepackage{amsmath,amssymb,mathtools}
\usepackage[scr=rsfs,cal=euler]{mathalfa}
\usepackage{xfrac}
\usepackage[unicode,hidelinks,bookmarks=false]{hyperref}
\newcommand{\ie}{i.e.\ }
\newcommand{\cf}{cf.\ }
\newcommand{\resp}{respectively\ }
\newcommand{\Subsection}{\subsection}
\providecommand{\nobreakdash}{\nobreak\hbox{-}}
\setlength{\emergencystretch}{2em}
\multlinegap0pt
\usepackage{amssymb}
\usepackage{float}
\usepackage{enumerate}
\newtheorem{theorem}{Theorem}
\newtheorem{claim}{Claim}
\title{Local measures of interval edge-uncolorability}
\author{Carl Johan Casselgren, Petros A. Petrosyan}
\date{Source: arXiv:2609.15873v1 (CC BY 4.0)}
\begin{document}
\maketitle

\noindent\emph{Source and licence.} Local measures of interval edge-uncolorability. Version arXiv:2609.15873v1 (CC BY 4.0), distributed by arXiv under the Creative Commons Attribution 4.0 International licence, http://creativecommons.org/licenses/by/4.0/, as recorded in the arXiv licence metadata for this exact version and shown on the abstract page https://arxiv.org/abs/2609.15873v1. Full licence notice: \texttt{licenses/arxiv-2609.15873-CC-BY-4.0.txt}.\par\smallskip

\noindent\textit{Editorial scope.} Counted unit: the article's theorem theorem (source0.tex lines 326--329). The article's complete original proof of it is delivered with it (source0.tex lines 330--545, one \texttt{proof} environment of 216 lines). No mathematics is written, repaired or re-derived: the statement and the proof are the article's own lines, in the article's own notation, and no retained line is substituted or re-flowed.  No further source-local context is needed: every cross-reference of every retained line resolves inside the two delivered spans.  Nothing was omitted from any retained span, so the package carries no omission record.  No retained line contains a citation, so the wrapper carries no bibliography and no bibliography entry.  No retained line contains a figure environment, an image-inclusion command or drawing code, so no image is delivered and the package declares no asset dependency.  Editorial formatting only, in addition: an AMS \texttt{amsart} preamble that restates the article's own declarations and macros used by the retained lines, namely the environments \texttt{theorem}, \texttt{claim} on separate counters, together with the standard packages those lines need.  The retained environments print in this document as Theorem 1, Claim 1 in order of appearance, whereas the article prints the same items with the numbering of its own full document.  The complete native source of the article as distributed in the arXiv e-print for this exact version is bundled unchanged as \texttt{sources/210457/source0.tex}.
\begin{theorem}
	Every graph with maximum degree at most $5$, where no two vertices of degree $3$ are
	adjacent has a weak near-interval coloring.
\end{theorem}

\begin{proof}
	Let $G$ be a graph as in the theorem;
	we color it as follows. Let $M$ be a maximum matching in the subgraph of $G$ induced
	by all vertices of degree $5$. Then $G-M$ satisfies that the vertices of degree $5$ induce
	an acyclic graph, so it has a proper $5$-edge coloring. We color the edges of $M$ by color $6$.
	This yields a proper $6$-edge coloring $\varphi$ of $G$, where only vertices of degree $5$ are incident with
	edges of color $6$.

	A vertex $v$ of $G$ where $\varphi$ is not weakly near-interval, we call {\em $\varphi$-problematic} (or just {\em problematic}).
	
	Straightforward case analysis yields that if $\varphi$ is problematic at a vertex $v$, then
\begin{itemize}

	\item $\varphi(v)=\{1,4\}$, $f(v)=\{1,5\}$, $f(v)=\{2,5\}$, if $v$ has degree $2$

	\item $\varphi(v)=\{1,4,5\}$, $f(v)=\{1,2,5\}$, if $v$ has degree $3$.

\end{itemize}

	An edge of color $1$ or $5$ that is incident with a problematic vertex is called {\em bad}.
	
	Using the coloring $\varphi$ as starting point, we shall prove that $G$ has a weak near-interval
	coloring by steps. We shall use colors $-1, \dots, 7$ in the process of constructing such a coloring.


	We first prove the following claim.

\begin{claim}
\label{cl:prob2}
	There is a proper edge coloring $g$ of $G$ with colors $0, \dots, 6$ such that
\begin{itemize}
	\item[(a)] no vertices of degree $2$ are problematic

	\item[(b)] colors $0$ and $6$ only appear on edges $e$ where one endpoint has degree at least $4$, and $g$
	is weakly near-interval at both endpoints of $e$,

	\item[(c)] no vertex is incident with exactly four edges that are colored $0,1,2,3$ or $3,4,5,6$.
\end{itemize}
\end{claim}

%	We first recolor (properly) some bad edges incident with vertices of degree 
%	$2$, using colors $0,\dots, 6$, in such a way that color 
%	$0$ or $6$ only appears at vertices
%	of degree at least $4$ that are adjacent to bad vertices, all recolored
%	edges are incident with vertices where the obtained coloring is weakly near-interval,
%	and no vertex  is incident with exactly four edges colored $0,1,2,3$, or $3,4,5,6$.

\begin{proof}	
	Let us assume that $g$ is a proper edge coloring, obtained from $\varphi$ 
	by recoloring some bad edges incident with vertices of degree $2$,
	with the fewest problematic vertices of degree $2$.
	Let us first note that if two problematic vertices are adjacent, then we can easily obtain a proper coloring
	with fewer problematic vertices, a contradiction. Thus we assume that this is not the case.

	Assuming that there is at least one problematic vertex of degree $2$ under $g$,
	we shall argue that there is a coloring $g'$ using colors $0, \dots, 6$ satisfying (a)-(c) and with
	fewer problematic vertices of degree $2$.

\bigskip

	{\bf Case 1.} {\em There is a problematic vertex $v$ which is incident with edges colored $1$ and $4$:}

	Suppose that $uv$ is colored $1$ under $f$, and that $vw$ is colored $4$. If color $2$ or $3$ does
	not appear at $u$, then we can recolor $uv$ by color $2$ or $3$ to obtain a coloring with fewer
	problematic vertices, so we assume that this is the case; analogously for $w$. Similary, color $5$
	must appear at $u$. If color $6$ does not appear at $u$, then we recolor $uv$ by color $6$.

	Otherwise, assume that color $6$ appears at $u$, so $g(u)=\{1,2,3,5,6\}$. If $5$ does not appear
	at $w$, then we recolor $uv$ by $4$ and $vw$ by $5$, so we may assume that this holds. Now, if
	$6$ does not appear at $w$, then we recolor $uv$ by $4$ and $vw$ by $6$, so we assume that
	$g(w) = \{2,3,4,5,6\}$. Hence, we can recolor $vw$ by the color $0$ to obtain a coloring
	with fewer problematic vertices, and where $v$ is not problematic.

	\bigskip

	{\bf Case 2.}	 {\em There is a problematic vertex $v$ which is incident with edges colored $2$ and $5$:}
	
	This case is completely analogous to the preceding case, where every color $c$ is replaced by $6-c$.

	\bigskip

	{\bf Case 3.}	 {\em There is a problematic vertex $v$ which is incident with edges colored $1$ and $5$:}
	
	Suppose $uv$ is colored $1$, and $vw$ is colored $5$. As in the preceding cases, we may assume that
	colors $3,4$ both appear at $u$, and colors $2,3$ both appear at $w$. Furthermore, we may
	assume that $4$ appears at $w$, and $2$ appears at $u$, because otherwise we may recolor $uv$
	or $vw$ and proceed as in Case 1 or Case 2. If $0$ does not appear at $w$, then we may recolor $vw$
	by this color, so we may assume that $0$ appears at $w$; similarly with color $6$ at $u$.
	Thus we may recolor $uv$ by $0$, and $vw$ by color $1$, to obtain a coloring with fewer
	problematic vertices.

	\bigskip
	
	This shows that there is a coloring $g'$ as required, which contradicts the choice of $g$. 
\end{proof}


	From the preceding claim we infer that
	there is a proper edge coloring $g$ of $G$ using colors $0,\dots, 6$ such that
\begin{itemize}

	\item $g(v)=\{1,4,5\}$ or $g(v)=\{1,2,5\}$, if $v$ is problematic under $g$,

	\item colors $0$ and $6$ only appear on edges $e$ where one endpoint has degree at least $4$, and $g$
	is weakly near-interval at both endpoints of $e$,

	\item no vertex is incident with exactly four edges that are colored $0,1,2,3$ or $3,4,5,6$.

\end{itemize}

	Next, we prove the following claim.


\begin{claim}
\label{cl:prob3}
	There is a proper edge coloring $f$ of $G$ with colors $-1, \dots, 7$ such that
\begin{itemize}
	\item[(i)] no vertices of $G$ are problematic

	\item[(ii)] colors $0$ and $6$ only appear on edges $e$ with one endpoint of degree at least $4$, 
	and $f$ is weakly near-interval at both endpoints of $e$, 

	\item[(iii)] no vertex is incident with exactly four edges that are colored $0,1,2,3$ or $3,4,5,6$.

	\item[(iv)] colors $-1$ and $7$ only  appear on edges that are not adjacent to bad edges.
\end{itemize}
\end{claim}
\begin{proof}
	Among the colorings satisfying (ii)-(iv), and with no problematic vertices of degree
	$2,4$ or $5$,
	we choose one $f$ with the minimum number
	of problematic vertices of degree $3$. Assuming that at least one vertex is problematic under $f$,
	we shall prove that there is a coloring $f'$ with fewer problematic vertices of degree $3$ 
	satisfying (ii)-(iv) (and with no problematic vertices of degree $2$, $4$, or $5$), 
	thereby deriving a contradiction.

%	Moreover, edges that are recolord $0$ or $6$ will not be adjacent to edges colored $1$ or $5$ that
%	are bad, i.e. incident with a problematic vertex.
	
	\bigskip

	{\bf Case 1.}	 {\em There is a problematic vertex $v$ with $f(v) =\{1,2, 5\}$:}

	Suppose that $uv$ is colored $1$, $vw$ is colored $2$, and $vx$ is colored $5$. As before,
	we may assume that $u,w$ and $x$ are all incident with an edge colored $3$. Moreover, color $4$
	appears at both $u$ and $x$.

	\bigskip

	{\bf Case 1.1.} {\em Color $1$ or $2$ appears at $x$:}

	Suppose that color $1$ or $2$ appears at $x$. If color $0$ does not appear at $x$, then we recolor
	$vx$ by $0$ and are done. So suppose that $0$ appears at $x$. If color $2$ appears at $x$, then
	we recolor $vx$ by $-1$. Otherwise, if $1$ appears at $x$, then we recolor $xv$ by $-1$ unless
	$x$ is adjacent to a problematic vertex $y$ via the edge colored $1$.
	
	So assume that there is such a problematic vertex $y$.
	If $f(y) = \{1,4,5\}$, then we recolor $xy$ by $2$ and $vx$ by $-1$, and are done.
	Otherwise if $f(y) = \{1,2,5\}$, then we set $v_1 = v$, $x_1 = x$, and consider a maximal
	$(1,5)$-colored path $P = v_1 x_1 v_2 x_2 \dots v_k x_k$ satisfying that
\begin{itemize}

	\item $v_1, \dots v_k,$ are problematic

	\item colors $f(v_i)=\{1,2,5\}$, $i=1,\dots, k-1$

	\item colors $f(v_k)=\{1,2,5\}$ or $f(v_k)=\{1,4,5\}$.

\end{itemize}	
	Note that we may assume that $\{3,4\} \subseteq f(x_i)$,  $i=1,\dots k$, 
	since otherwise we can recolor $v_i x_i$
	and obtain a coloring with fewer problematic vertices.

	Now we recolor the edges $v_1x_1, \dots v_k x_k$ as follows:

\begin{itemize}

	\item If $0$ appears at $x_i$, then we recolor $v_i x_i$ by color $-1$, otherwise 
	we recolor it by $0$, for $i=1,\dots, k-1$.
	
	\item If $1,4,5$ appear at $v_k$ and $2$ does not appear at $x_{k-1}$, then we recolor $x_{k-1}v_k$
 	by $2$. On the other
	hand, if $2$ appears at $x_{k-1}$, then we recolor it $6$.


	\item If $1,2,5$ appear at $v_k$, then since $x_5$ has degree at least $4$, and 
	$f(x_k) \neq \{3,4,5,6\}$, $1$ or $2$ must appear at $x_5$.
	If $0$ does not appear at $x_k$, then we recolor $v_kx_k$ by $0$;
	otherwise if $0$ does appear at $x_k$, then we recolor $x_k v_k$ by $-1$.
\end{itemize}

	It is readily verified that this yields a required coloring $f'$, which contradicts the choice of $f$.

	\bigskip

	{\bf Case 1.2.} {\em Neither of color $1$ and $2$ appear at $x$:}

	Note that the conditions (i)-(iv) imply that $0$ does not appear at $x$. 
	Moreover, since $x$ does not have degree $3$, $6$ must appear at $x$.
	By the same conditions, $7$ cannot appear at $x$, so
	$x$ has degree $4$ and colors $3,4,5,6$ are present at $x$.
	However, this contradicts the condition (iii) above.
	We conclude that Case 1.2 is not possible.


	\bigskip

	{\bf Case 2.}	 {\em There is a problematic vertex $v$ with $f(v) = \{1,4,5\}$:}

	This case is completely analogous to Case 1 by replacing every color $c$ in that case by $6-c$.
\end{proof}	
\bigskip

	By the preceding claim we may conclude that there is a coloring satisfying (i)-(iv) with no problematic
	vertices of degree $3$. This is the required weak near-interval coloring.
\end{proof}

\end{document}
